\section{Spanning tree}
\label{sec:st}

Let $H_{n, k}$ denote a directed graph with $V(H_{n,k}) = \{v_{c,r} : c \in [n], r \in [2k]\}$ and
\begin{align*}
    E(H_{n, k}) = \{(v_{c_1, r_1}, v_{c_2, r_2}) &\mid
    c_1, c_2 \in [n], r_1, r_2 \in [2k], |c_1 - c_2| = 1 \} \ \cup \\
    \{(v_{c, r_1}, v_{c, r_2}) &\mid
    c \in \{1, n\}, r_1, r_2 \in [2k], r_1 \neq r_2 \}.
\end{align*}

\begin{figure}[h]
    \centering
    \label{fig:dst}
    \begin{tikzpicture}
		\node [style=node] (0) at (-3.5, 1.5) [label={$v_{1,1}$}] {};
		\node [style=node] (1) at (-3.5, 0.5) [label={$v_{1,2}$}] {};
		\node [style=node] (2) at (-3.5, -0.5) [label={$v_{1,3}$}] {};
		\node [style=node] (3) at (-3.5, -1.5) [label={$v_{1,4}$}] {};
		\node [style=node] (4) at (-2, 1.5) {};
		\node [style=node] (5) at (-2, 0.5) {};
		\node [style=node] (6) at (-2, -0.5) {};
		\node [style=node] (7) at (-2, -1.5) {};
		\node [style=none] (16) at (-0.5, 1.5) {};
		\node [style=none] (17) at (-0.5, 0.5) {};
		\node [style=none] (18) at (-0.5, -0.5) {};
		\node [style=none] (19) at (-0.5, -1.5) {};
		\node [style=none] (20) at (-0.5, -1.25) {};
		\node [style=none] (21) at (-0.5, -0.75) {};
		\node [style=none] (22) at (-0.5, -0.25) {};
		\node [style=none] (23) at (-0.5, 0.25) {};
		\node [style=none] (24) at (-0.5, 0.75) {};
		\node [style=none] (25) at (-0.5, 1.25) {};
		\node [style=node] (26) at (3.5, 1.5) [label={$v_{n,1}$}] {};
		\node [style=node] (27) at (3.5, 0.5) [label={$v_{n,2}$}] {};
		\node [style=node] (28) at (3.5, -0.5) [label={$v_{n,3}$}] {};
		\node [style=node] (29) at (3.5, -1.5) [label={$v_{n,4}$}] {};
		\node [style=node] (30) at (2, 1.5) {};
		\node [style=node] (31) at (2, 0.5) {};
		\node [style=node] (32) at (2, -0.5) {};
		\node [style=node] (33) at (2, -1.5) {};
		\node [style=none] (34) at (0.5, 1.5) {};
		\node [style=none] (35) at (0.5, 0.5) {};
		\node [style=none] (36) at (0.5, -0.5) {};
		\node [style=none] (37) at (0.5, -1.5) {};
		\node [style=none] (38) at (0.5, -1.25) {};
		\node [style=none] (39) at (0.5, -0.75) {};
		\node [style=none] (40) at (0.5, -0.25) {};
		\node [style=none] (41) at (0.5, 0.25) {};
		\node [style=none] (42) at (0.5, 0.75) {};
		\node [style=none] (43) at (0.5, 1.25) {};
		\node [style=none] (44) at (0, 0) {$\dots$};
        
		\draw [style=bi-directional] (0) to (4);
		\draw [style=bi-directional] (0) to (5);
		\draw [style=bi-directional] (0) to (6);
		\draw [style=bi-directional] (0) to (7);
		\draw [style=bi-directional] (1) to (4);
		\draw [style=bi-directional] (1) to (5);
		\draw [style=bi-directional] (1) to (6);
		\draw [style=bi-directional] (1) to (7);
		\draw [style=bi-directional] (2) to (4);
		\draw [style=bi-directional] (2) to (5);
		\draw [style=bi-directional] (2) to (6);
		\draw [style=bi-directional] (2) to (7);
		\draw [style=bi-directional] (3) to (4);
		\draw [style=bi-directional] (3) to (5);
		\draw [style=bi-directional] (3) to (6);
		\draw [style=bi-directional] (3) to (7);
		\draw [style=bi-directional, bend right=90, looseness=1.25] (0) to (3);
		\draw [style=bi-directional, bend right=75, looseness=1.25] (1) to (2);
		\draw [style=bi-directional, bend right=75, looseness=1.25] (0) to (1);
		\draw [style=bi-directional, bend right=75, looseness=1.25] (2) to (3);
		\draw [style=bi-directional, bend right=75, looseness=1.25] (0) to (2);
		\draw [style=bi-directional, bend right=75, looseness=1.25] (1) to (3);
		\draw [style=directed] (16.center) to (4);
		\draw [style=directed] (17.center) to (5);
		\draw [style=directed] (18.center) to (6);
		\draw [style=directed] (19.center) to (7);
		\draw [style=directed] (20.center) to (6);
		\draw [style=directed] (20.center) to (5);
		\draw [style=directed] (20.center) to (4);
		\draw [style=directed] (21.center) to (7);
		\draw [style=directed] (22.center) to (5);
		\draw [style=directed] (22.center) to (4);
		\draw [style=directed] (23.center) to (7);
		\draw [style=directed] (23.center) to (6);
		\draw [style=directed] (24.center) to (4);
		\draw [style=directed] (25.center) to (7);
		\draw [style=directed] (25.center) to (6);
		\draw [style=directed] (25.center) to (5);
		\draw [style=bi-directional] (26) to (30);
		\draw [style=bi-directional] (26) to (31);
		\draw [style=bi-directional] (26) to (32);
		\draw [style=bi-directional] (26) to (33);
		\draw [style=bi-directional] (27) to (30);
		\draw [style=bi-directional] (27) to (31);
		\draw [style=bi-directional] (27) to (32);
		\draw [style=bi-directional] (27) to (33);
		\draw [style=bi-directional] (28) to (30);
		\draw [style=bi-directional] (28) to (31);
		\draw [style=bi-directional] (28) to (32);
		\draw [style=bi-directional] (28) to (33);
		\draw [style=bi-directional] (29) to (30);
		\draw [style=bi-directional] (29) to (31);
		\draw [style=bi-directional] (29) to (32);
		\draw [style=bi-directional] (29) to (33);
		\draw [style=bi-directional, bend left=90, looseness=1.25] (26) to (29);
		\draw [style=bi-directional, bend left=75, looseness=1.25] (27) to (28);
		\draw [style=bi-directional, bend left=75, looseness=1.25] (26) to (27);
		\draw [style=bi-directional, bend left=75, looseness=1.25] (28) to (29);
		\draw [style=bi-directional, bend left=75, looseness=1.25] (26) to (28);
		\draw [style=bi-directional, bend left=75, looseness=1.25] (27) to (29);
		\draw [style=directed] (34.center) to (30);
		\draw [style=directed] (35.center) to (31);
		\draw [style=directed] (36.center) to (32);
		\draw [style=directed] (37.center) to (33);
		\draw [style=directed] (38.center) to (32);
		\draw [style=directed] (38.center) to (31);
		\draw [style=directed] (38.center) to (30);
		\draw [style=directed] (39.center) to (33);
		\draw [style=directed] (40.center) to (31);
		\draw [style=directed] (40.center) to (30);
		\draw [style=directed] (41.center) to (33);
		\draw [style=directed] (41.center) to (32);
		\draw [style=directed] (42.center) to (30);
		\draw [style=directed] (43.center) to (33);
		\draw [style=directed] (43.center) to (32);
		\draw [style=directed] (43.center) to (31);
\end{tikzpicture}

    \caption{$H_{n,2}$}
\end{figure}

It is easy to see that the pathwidth of $H_{n, k}$ is at most $4k$. Combining the following lemma with \Cref{lem:raz_spieker} immediately gives \Cref{thm:main_dst}.

\begin{lemma}\label{lem:tc_dst}
    For any $k \geq 1$ and $n > 12k^2 \ln k + 3$, any tropical circuit calculating $DST_{H_{n, k}}$ uses at least
    %$2^{\Omega(k \log \log k)}$ gates.
    $C_k / (2k^2 \ln k)$
    gates.
\end{lemma}

The rest of this section is dedicated to proving this lemma. The majority of the proof follows the same line of argumentation used to prove \Cref{lem:main_dtsp}.
For the rest of the section, we fix suitable $n, k$ and put $G := H_{n, k}$. Let $V_i = \{v_{i, r} : r \in [2k]\}$. For $i \in [n - 1]$, let $G^i = G[V_i \cup V_{i + 1}] - (E(G[V_i]) \cup E(G[V_{i+1}]))$ and let $G^0 = G[V_1], G^n = G[V_n]$. Note that $G^i$ is complete bipartite for $1 \leq i \leq n - 1$ and complete for $i \in \{0, n\}$. We denote $E_i = E(G^i)$ for all $i \in \{0, \dots, n\}$ and put $X_i = \{x_e \colon e \in E_i\}$.

For every perfect matching $M$ of $G^i$ for $i \in \{1, \dots, n-1\}$, we will identify it with a permutation $\rho \in \sym_{2k}$ defined as
    $$
        \rho(s) = t \quad \text{iff} \quad (v_{i, s}, v_{i + 1, t}) \in M \ \text{or}\ (v_{i+1, t}, v_{i,s}) \in M.
    $$
    Similarly, for $i \in \{0, n\}$ we identify any perfect matching $M$ of $G^i$ with $\rho \in \sym_{2k}^2$ such that
    $$
        \rho(s) = t \quad \text{iff} \quad (v_{i, s}, v_{i, t}) \in M \ \text{or}\ (v_{i, t}, v_{i, s}) \in M.
    $$
    
    
    We will say that a Hamiltonian cycle of $G$ is \emph{nice} if its intersection with each $E_i$ is a perfect matching in $G^i$ (ignoring edge directions). We will say that a Hamiltonian path of $G$ is \emph{nice} if it is a subgraph of a nice Hamiltonian cycle and its first and last vertices belong to $V_1$. Every nice path is a subgraph of exactly one nice cycle, and conversely, every nice cycle contains exactly $k$ nice paths as a subgraph.

    We define the content of a nice Hamiltonian cycle $H$ as the unique sequence of permutations $\rho_0, \dots, \rho_n \in \sym_{2k}$ such that $E(H) \cap E_i = \rho_i$ for all $0 \leq i \leq n$.
    For the sake of clarity, for such a sequence, we define the following notation.
    $$
        \overleftarrow{\rho_i} := \rho_i \rho_{i - 1} \dots \rho_2 \rho_1
    $$
    $$
        \overrightarrow{\rho_i} := \rho_{n - 1} \rho_{n - 2} \dots \rho_{i + 1} \rho_{i}
    $$
    We assume $\overleftarrow{\rho_0} = \overrightarrow{\rho_n} = \textbf{id}_{\sym_{2k}}$. 
    Note that as described in \Cref{sec:mm*},
    \begin{equation}\label{eq:dst_nice_cycles}
        \left(
            \overleftarrow{\rho_{i - 1}} \cdot
            \rho_0 \cdot
            \overleftarrow{\rho_{i - 1}}^{-1}
        \right)
        \cdot
        \left(
            \overrightarrow{\rho_i}^{-1} \cdot
            \rho_n \cdot
            \overrightarrow{\rho_i}
        \right)
        \in \sym_{2k}^k
    \end{equation}
    for any $i \in [n]$.
    Conversely, for every sequence of permutations satisfying the above (for arbitrarily chosen $i$), there are exactly two nice cycles $H$ such that their content is $\rho_0, \dots, \rho_n$. Both cycles differ only by their direction. In particular, if \Cref{eq:dst_nice_cycles} holds for one such $i$, then it holds for every $i \in [n]$.

    \begin{claim}
        There are $((2k)!)^{n - 1} \cdot (2k - 1)!$ sequences $(\rho_j)_{j=0}^n$ satisfying (1).
    \end{claim}

    \begin{proof}
        We fix $\rho_j$ arbitrarily for all $j \neq i$; there are $((2k - 1)!!)^2 \cdot ((2k)!)^{n - 2}$ ways to do that. Let $\bar\rho = \left(\overrightarrow{\rho_i}^{-1} \cdot \rho_n \cdot \overrightarrow{\rho_i}\right)$, hence we have $\bar\rho = \rho_i^{-1} \tilde\rho \rho_i$, where $\tilde\rho \in \sym_{2k}^2$ is already fixed. By \Cref{prop:sym_k}, the number of values $\bar\rho$ can take for \Cref{eq:dst_nice_cycles} to hold is $(2k - 2)!!$. By the same proposition, the number of $\rho_i \in \sym_{2k}^2$ given $\bar\rho$ is $(2k)!!$. The total number of ways to fix all $\rho_j$ is therefore $((2k - 1)!!)^2 \cdot ((2k)!)^{n - 2} \cdot (2k)!! \cdot (2k - 2)!! = ((2k)!)^{n - 1} \cdot (2k - 1)!$.
    \end{proof}
    
    In particular, the number of nice cycles in $G$ is $2 \cdot ((2k)!)^{n-1} \cdot (2k - 1)!$. By a similar argument, for any set of indices $1 \leq i_1, \dots, i_\ell \leq n - 1$ of size at most $n - 2$, if we fix $\rho_{i_1}, \dots, \rho_{i_\ell}$ together with $\rho_0$ and $\rho_n$, then there are exactly
    $
        ((2k)!)^{n - 2 - \ell} \cdot (2k)!! \cdot (2k - 2)!!
    $
    ways to fix the rest of the permutations for (1) to hold.
    \newline




\begin{proof}[Proof (\Cref{lem:tc_dst})]

    We say that a rectangle $g \cdot h \subseteq DST_G$ covers a nice Hamiltonian cycle $H$ if it contains a monomial encoding a nice Hamiltonian path which is a subgraph of $H$. Since all nice Hamiltonian paths are spanning out-trees as well, any rectangle decomposition of the polynomial calculated by our circuit must cover all nice cycles. We define rectangles being balanced as in the statement of \Cref{lem:tc_decomposition} with respect to the set $X = \{ x_{(v_{c, 1}, v_{c + 1, 1})} \mid c \in [n - 1] \}$.
    
    \begin{claim}\label{clm:dst_rect_bound}
        Let $g \cdot h$ be any balanced rectangle in $DST_G$. Then, the number of nice Hamiltonian cycles covered by $g \cdot h$ is at most
        $$
            ((2k)!)^n \cdot C_k^{-1} \cdot 2k \ln k.
        $$
    \end{claim}

    For any tropical circuit calculating $DST_G$ using $\tau$ gates, \Cref{lem:tc_decomposition} implies that $DST_G$ can be covered by a union of $\tau$ balanced rectangles. Thus, in order to cover every nice cycle, at least one balanced rectangle must cover at least $\frac{2 ((2k)!)^{n - 1} (2k - 1)!}{\tau}$ of those, hence assuming \Cref{clm:dst_rect_bound},
    $$
        \tau \geq C_k / (2k^2 \ln k).
    $$    
    
    The rest of the section is dedicated to proving \Cref{clm:dst_rect_bound}.
    Fix any balanced rectangle $g \cdot h$. The supports of $g$ and $h$ are disjoint. We will say that $E_i$ is monochromatic w.r.t. $g \cdot h$ if either $X_i \cap \sup(g)$ or $X_i \cap \sup(h)$ is empty.

    \begin{claim}\label{clm:dst_rect_colorful}
        Pick $1 \leq i \leq n - 1$. If $E_i$ is not monochromatic w.r.t. $g \cdot h$, then there exists an edge $(u, v) \in E_i$ such that $x_{u, v} \not\in \sup(g) \cup \sup(h)$.
    \end{claim}

    \begin{proof}
        Assume by contradiction that $x_{u, v} \in \sup(g)$ or $x_{u, v} \in \sup(h)$ for every $(u, v) \in E_i$.
        Since all monomials in $g \cdot h$ represent out-trees, for any pair of edges of the form $(u, v), (w, v) \in E_i$, either both $x_{u, v}, x_{w, v}$ belong to $\sup(g)$, or both belong to $\sup(h)$. Similarly, for any $(u, v) \in E_i$ both $x_{u, v}, x_{v, u}$ must also belong to the same support.
        
        By our assumption, we have two edges $(p, q), (s, t) \in E_i$ such that $x_{p, q} \in \sup(g)$ and $x_{s, t} \in \sup(h)$ (and hence $x_{t, s} \in \sup(h)$ as well). Either $(q, s)$ or $(q, t)$ belong to $E_i$. In the former case, we have $x_{q, s} \in \sup(h)$ and $x_{s, q} \in \sup(g)$. In the latter, $x_{q, t} \in \sup(h)$ and $x_{t, q} \in \sup(g)$. Both cases arrive at a contradiction, which proves the claim.
    \end{proof}

    
    Let $i_1 < \dots < i_\ell$ denote the indices in range $[n - 1]$ of sets $E_i$ which are not monochromatic. We consider two cases depending on whether $\ell > 4k^2 \ln k$ or not.

    Assume $\ell > 4k^2 \ln k$. By sacrificing at most half of indices $i_j$, we can assume that they are non-adjacent. That is, we assume that $\ell > 2k^2 \ln k$ and that $i_{j + 1} - i_j > 1$ for all $j \in [\ell - 1]$. For each $i \in \{i_j\}_{j=1}^\ell$, let $(p_i, q_i)$ denote the edge of $E_i$ given by \Cref{clm:dst_rect_colorful}.
    Consider the following randomized procedure to generate a nice directed Hamiltonian cycle $H$ of $G$.
    \begin{itemize}
        \item Select $\rho_0 \in \sym_{2k}^2$ uniformly randomly.
        \item Select a subset of edges $F_0 \subseteq E_0$ by matching vertices of $V_1$ according to $\rho_0$ and picking the edge directions uniformly and independently. Let $V_1^{tail}$ and $V_1^{head}$ denote the vertices of $V_1$ composed of respectively tails and heads of the edges of $F_0$.
        \item For $i \in \{1, \dots, n - 1\}$ do the following. Pick $\rho_i \in \sym_{2k}$ uniformly randomly and select $F_i \subseteq E_i$ by matching vertices of $V(G^i)$ according to $\rho_i$ and picking the edge directions so that they agree with edges of $F_{i - 1}$ (i.e., if we match $v_{i, r_1}$ with $v_{i + 1, r_2}$ and $v_{i, r_1} \in V_i^{tail}$, then we add $(v_{i + 1, r_2}, v_{i, r_1})$ to $F_i$, and $(v_{i, r_1}, v_{i + 1, r_2})$ otherwise). Define $V_{i + 1}^{tail}, V_{i + 1}^{head}$ analogously.
        \item The union $F_0 \cup \dots \cup F_{n - 1}$ at this point forms $k$ disjoint directed paths whose endpoints are exactly $V_n$. There are exactly $(k - 1)!$ ways to pick $F_n$ in a way that connects these paths into a Hamiltonian cycle. We pick one of such ways uniformly randomly.
    \end{itemize}
    It is easy to see that this procedure generates all nice Hamiltonian cycles of $G$ with uniform probability distribution. Now, we can bound the probability that a cycle generated this way does not contain any of the forbidden edges $(p_i, q_i)$ given by \Cref{clm:dst_rect_colorful}. Note that only such cycles can be covered by $g \cdot h$.

    The steps during which the number of feasible choices could get restricted is when selecting $F_{i_j}$ for some $j \in [\ell]$. Assume w.l.o.g. $p_{i_j} \in V_{i_j}, q_{i_j} \in V_{i_j + 1}$, the other case is symmetric. If $p_{i_j} \in V_{i_j}^{tail}$, then every choice of $F_{i_j}$ will be feasible, as an edge between $p_{i_j}, q_{i_j}$ potentially added will always be directed from $q_{i_j}$ to $p_{i_j}$. In case $p_{i_j} \in V_{i_j}^{head}$, there is exactly $\frac{1}{k}$ chance of adding the forbidden edge to $F_{i_j}$, as $\rho_{i_j}$ is selected uniformly from whole $\sym_{2k}$.

    The crucial observation is that the probability that $p_{i_j}$ is in $V_{i_j}^{head}$ at this point, assuming so far we did not pick any of the forbidden edges, is exactly $1/2$, independent of previous choices. This follows from the fact that $\rho_{i_j - 1}$ was chosen uniformly from $\sym_{2k}$ and due to our assumption that $i_j > i_{j - 1} + 1$. The partition of $V_{i_j - 1}$ into $V_{i_j - 1}^{head}$ and $V_{i_j - 1}^{tail}$ can have arbitrary distribution, however, both sets are always of size $k$, hence the probability of $p_{i_j}$ getting matched by $\rho_{i_j - 1}$ with a vertex from $V_{i_j - 1}^{tail}$ is exactly $\frac{1}{2}$.

    This means, that the probability of a nice Hamiltonian cycle containing none of the forbidden edges is at most
    $$
    \left( 1 - \frac{1}{2k} \right)^\ell < e^{-k \ln k} \leq C_k^{-1} \leq C_k^{-1} \cdot 2k^2 \ln k,
    $$
    hence the number of such cycles is at most
    $$
    \left( 2 ((2k)!)^{n-1} (2k - 1)! \right) 
        \cdot \left( C_k^{-1} \cdot 2k^2 \ln k \right) =
    ((2k)!)^n \cdot C_k^{-1} \cdot 2k \ln k,
    $$
    which finishes the proof in case $\ell > 4k^2 \ln k$.
    \newline
    

    Now, assume $\ell < 4k^2 \ln k$. Thus, at least $n - 4k^2 \ln k \geq \frac{2n}{3} + 1$ of $E_i$ are monochromatic. Since $g \cdot h$ is balanced, we have two indices $i_g, i_h \in [n - 1]$ s.t. $E_{i_g}, E_{i_h}$ are both monochromatic, and both $X_{i_g} \cap \sup(h)$ and  $X_{i_h} \cap \sup(g)$ are empty. W.l.o.g. we can assume that $i_g < i_h$. Let $\bar{E} = E(G) - (E_{i_g} \cup E_{i_h})$.

    For a monomial $g' \in g$ (and analogously for $h' \in h$), we define its type as the set
    $$
        \lambda_{g'} := \{ e \in \bar{E} : x_e \in \sup(g') \}.
    $$
    Fix an arbitrary pair of types $\lambda_{g^*}, \lambda_{h^*}$ and let $g^*$ (resp. $h^*$) denote the set of all monomials of $g$ (resp. $h$) with that type. By the definition, every spanning out-tree corresponding to some monomial of $g^* \cdot h^*$ has the same intersection with $\bar{E}$. If $g^* \cdot h^*$ covers at least one nice cycle, this intersection must induce a perfect matching in $G^i$ for each $i \in [n] - \{i_g, i_h\}$ and a matching of size $k - 1$ in $G^0$.
    Let $(\rho_i)_{i \in [n] - \{i_g, i_h\}}$ denote the set $(\lambda_{g^*} \cup \lambda_{h^*}) \cap E_i$. Let $\rho_0$ denote the unique permutation of $\sym_{2k}^2$ containing $(\lambda_{g^*} \cup \lambda_{h^*}) \cap E_0$ as a subgraph.
    
    Let $\bar\rho_l = \overleftarrow{\rho_{i_g - 1}}$, $\bar\rho_r = \overrightarrow{\rho_{i_h+1}}$ and $\bar\rho_m = \rho_{i_h - 1}\rho_{i_h - 2}\dots\rho_{i_g + 2}\rho_{i_g + 1}$. Let $P_g$ (resp. $P_h$) denote the set of perfect matchings of $G^{i_g}$ (resp. $G^{i_h})$ induced by the monomials in $g^*$ (resp. $h^*$). Finally, put
    $$
        P_g^* = \left\{
            \rho \cdot
            \bar\rho_l \cdot 
            \rho_0 \cdot 
            \bar\rho_l^{-1} \cdot 
            \rho^{-1}
            \mid \rho \in P_g
        \right\} \ \text{and}
    $$
    $$
        P_h^* = \left\{
            \bar\rho_m^{-1} \cdot
            \rho^{-1} \cdot
            \bar\rho_r^{-1} \cdot 
            \rho_n \cdot 
            \bar\rho_r \cdot 
            \rho \cdot
            \bar\rho_m
            \mid \rho \in P_h
        \right\}.
    $$
    By \Cref{prop:sym_k}, we have $|P_g| \leq (2k)!! \cdot |P_g^*|$ and $|P_h| \leq (2k)!! \cdot |P_h^*|$.
    By \Cref{eq:dst_nice_cycles}, every pair $\rho_g \in P_g^*, \rho_h \in P_h^*$ satisfies $\rho_h^* \rho_g^* \in \sym_{2k}^k$, hence applying \Cref{lem:mm*_size_bound} gives us
    $$
        2 \cdot |P_g| \cdot |P_h| \leq
        2 ((2k)!!)^2 \cdot |P_g^*| \cdot |P_h^*| \leq
        2 ((2k)!!)^2 (2k - 1)! \cdot C_k^{-1} \cdot 2k \ln k.
    $$
    Finally, note that the quantity $2 \cdot |P_g| \cdot |P_h|$ bounds from above the number of nice cycles covered by $g^* \cdot h^*$.

    The number of pairs of types $\lambda_{g^*}, \lambda_{h^*}$ which cover at least one nice cycle is bounded by the number of different projections of a nice path onto $\bar{E}$, hence by
    $$
    k ((2k - 1)!!)^2 ((2k)!)^{n - 3}.
    $$
    Therefore, the total number of nice cycles covered by $g \cdot h$ is at most
    $$
    \left( k ((2k - 1)!!)^2 ((2k)!)^{n - 3} \right) \cdot
    \left( 2 ((2k)!!)^2 (2k - 1)! \cdot C_k^{-1} \cdot 2k \ln k \right) =
    $$
    $$
    = ((2k)!)^n \cdot C_k^{-1} \cdot 2k \ln k,
    $$
    which finishes the proof of the claim.

\end{proof}